\documentclass[UTF8,a4paper,11pt]{ctexart}
\usepackage[margin=2.5cm]{geometry}
\usepackage{array,longtable,xcolor,url,amssymb}

\definecolor{PendingColor}{RGB}{180,82,0}
\renewcommand{\arraystretch}{1.35}
\setlength{\parindent}{2em}
\setlength{\parskip}{0.25em}

\newcommand{\record}[9]{%
  \section*{#1}
  \begin{longtable}{>{\bfseries}p{0.22\textwidth}p{0.70\textwidth}}
  工具与模型 & #2 \\
  使用目的 & #3 \\
  主要提示方式 & #4 \\
  使用过程 & #5 \\
  采纳情况 & #6 \\
  人工修改 & #7 \\
  人工核验 & #8 \\
  关联成果位置 & #9 \\
  \end{longtable}
}

\begin{document}

\begin{center}
  {\LARGE\bfseries AI工具使用详情}
\end{center}

\noindent\textbf{填写状态：}\textcolor{PendingColor}{\textbf{工作稿，提交前须由参赛队逐项核验并签认。}}

本队在竞赛过程中使用了 Codex。核心建模选择、正式测试操作和最终结论由参赛队负责；下列 AI 参与成果在提交前必须由队员对照题面、代码、日志和数值结果逐项人工审查。本文件不记录账号、密码、令牌或身份信息。

\record{记录 01：仓库与规则准备}
{Codex，GPT-6；产品精确构建版本未获取。}
{整理 2026 竞赛规则、论文格式、AI 使用规定、项目模板、检查器及 Python、MATLAB 环境。}
{要求保留已有修改，按官网材料更新规则与模板，检查环境并初始化正式比赛目录。}
{读取官方来源记录，建立规则映射、检查器、论文和 AI 详情模板；运行 Python、MATLAB、LaTeX 与跨语言验证。}
{仓库结构、规则记录、模板和检查器已作为比赛工作目录使用。}
{队伍须在提交前复核规则日期、页面限制、材料命名和截止时间；目前未记录逐项修改人。}
{自动测试已运行；官方规则原文仍须由参赛队再次逐条核对。}
{\path{rules/}、\path{templates/}、\path{checkers/}、\path{environment/compatibility.md}。}

\record{记录 02：题面读取与任务拆解}
{Codex，GPT-6。}
{读取 B 题题面和附件，统一坐标、角度、时间、通信协议及四问交付物。}
{用户提供 B 题 PDF、附件压缩包和模拟器目录，并确认选择 B 题。}
{提取并渲染检查 4 页题面，读取附件中的干扰源、测向机和接口说明，形成任务表、协议说明与风险清单。}
{题面条件、四问依赖关系和接口状态机已进入项目工作稿。}
{参赛队须核对全部单位、接口字段和题面页码；目前未记录逐项修改人。}
{题面 PDF 与提取文本已交叉核对；模型含义仍须由队员独立复述确认。}
{B题任务分解.md、模拟器接入说明.md、\path{sources.yaml}。}

\record{记录 03：几何模型、代码与测试}
{Codex，GPT-6 与 GPT-5；产品精确构建版本未获取。}
{辅助推导示向楔形、定位区域直径、最小包围圆、第二检测点候选域、600 米覆盖证明及程序实现。}
{要求遵循仓库工作流，采用测试先行，所有结论必须可复算并登记结果清单。}
{先写失败测试，再实现几何函数和确定性分析脚本；构造等边三角形反例，推导三圆盘交集和对称第二检测点，并进行旋转不变性、边界与网格收敛测试。}
{几何代码、测试和 JSON 结果进入工作稿；主要数值为 $D=40$ 米、最小包围圆半径 $23.094$ 米及第二检测点局部坐标 $(761.429,\pm648.248)$ 米。}
{队员应独立复算半平面符号、三圆盘结论和推荐点公式，并记录实际修改。}
{自动测试与数值收敛已完成；队员人工推导签认尚未完成。}
{\path{python/b_geometry.py}、\path{python/q1_q2_analysis.py}、\path{output/q1-q2-geometry-20260911.json}。}

\record{记录 04：模拟器演练与策略迭代}
{Codex，GPT-6。}
{实现通信客户端、覆盖搜索和定位清除策略，运行合成模拟，并在用户确认演练页面后执行官方演练。}
{用户要求使用模拟测试并允许使用比赛文件；每次演练由用户确认模拟器状态。}
{先完成合成回归，再形成 v1、最近未完成点续扫 v2 和清除 16 个后停止 v3；问题三、四各完成三轮官方演练并保存动作日志和模拟器原始文件。未启动正式测试。}
{v3 已冻结为正式候选；六轮演练累计清除 $82/82$ 个源。}
{队伍须确认正式测试前继续采用冻结版，并核对每轮正式测试题号和轮次。}
{原始日志、结果文件和 SHA-256 已自动核对；正式测试仍需逐轮人工授权。}
{\path{python/b_client.py}、\path{python/b_strategy.py}、\path{python/b_simulation.py}、\path{logs/official-practice/}。}

\record{记录 05：图表和论文草稿}
{Codex，GPT-5。}
{根据几何 JSON 和官方演练结果制作图表，并辅助形成问题一至四论文工作稿。}
{要求图表不改变数据和结论；演练数据不得冒充正式成绩。}
{生成几何反例、第二检测点候选域和演练汇总图；编写模型、结果、检验、结论与 AI 声明，正式结果表保留“尚未执行”。}
{图表和论文进入工作稿。}
{队伍须逐段修改措辞，核对公式、图表和数值，并记录人员与内容。}
{自动来源核对、单元测试、论文编译和渲染检查已完成；队员逐页签认尚未完成。}
{\path{paper/}、\path{python/b_figures.py}、\path{figures/final/}。}

\record{记录 06：往年展示论文结构与排版分析}
{Codex，GPT-6。}
{提取往年官方展示论文的结构、图表和排版特征，建立当前工作稿的改进基准。}
{用户要求阅读往年获奖论文并学习其结构、论文格式和排版。}
{逐页审读 2025 年 B060、B157 和抽查 2023 年 A0175、C050 官方展示页，共处理 165 页并记录可观察证据。}
{形成显式问题分析、按问串联、图表承担证据和格式服从 2026 规则的写作基准。}
{队伍须决定最终章节顺序和篇幅，不能把展示特征解释为获奖因果。}
{页面、页眉和章节位置已核对；最终版式仍须由队员逐页审阅。}
{\path{knowledge/award-paper-structure-layout.md}、\path{workflows/writing-and-figures.md}。}

\record{记录 07：论文证据链与算法图补强}
{Codex，GPT-6。}
{使问题分析、有限算法、覆盖证明和文献依据可在正文中直接核查。}
{用户要求继续通信、几何和论文工作，并明确不启动正式测试。}
{保留冻结 v3 和已有数值，新增四问问题分析、有限控制流程图、600 米覆盖与定向半平面证明图，并补充最小包围圆、计算几何和规划算法文献；随后执行本地测试、合成回归、论文编译和逐页渲染。}
{新增正文和图表进入论文工作稿；正式结果表继续保留“尚未执行”。}
{队员须核对八次上限、$19.5$ 米阈值、61 个覆盖点、16 源退出条件和定向源证明。}
{自动验证结果见兼容性记录；参赛队逐页签认尚未完成。}
{\path{paper/sections/analysis.tex}、\path{paper/sections/model.tex}、\path{figures/final/q34-*}。}

\section*{最终人工复核（尚待完成）}

\begin{itemize}
  \item[$\square$] 每位队员能独立解释四问模型、公式、策略和正式测试限制；
  \item[$\square$] 正式测试全部完成后，将六轮结果与原始日志逐项核对，且演练数据不冒充正式数据；
  \item[$\square$] 已记录队员对 AI 草稿的实际采纳、修改和核验；
  \item[$\square$] 本文件不含账号、密码、令牌或个人身份信息；
  \item[$\square$] 最终版已导出为“AI工具使用详情.pdf”，且与论文声明一致。
\end{itemize}

\end{document}
